\vspace{0.5cm}
\section{Progettazione Logica}

In questa sezione viene descritto il processo di \textbf{reificazione} dello schema E-R presentato in precedenza.
L'obiettivo è rappresentare i dati in modo preciso ed efficiente, analizzando le ridondanze, eliminando le generalizzazioni ed infine illustrando graficamente il diagramma risultante, fino a ottenere una struttura complessiva ottimizzata per la traduzione in linguaggio SQL.

\vspace{0.5cm}
\subsection{Analisi delle Ridondanze}

Come evidenziato nello schema E-R in \textbf{Figura 1}, l'attributo \textit{Totale\_Vendite} effettuate da un venditore risulta essere ridondante. Tale informazione, infatti, può essere ottenuta direttamente contando le occorrenze del venditore all'interno della relazione Acquisto.

L'attributo viene attualmente aggiornato ad ogni nuova vendita (circa 2/3 al giorno) e consultato alla fine di ogni mese per monitorare l'andamento dei venditori dipendenti.

\vspace{0.5cm}

\renewcommand{\arraystretch}{1.5}
\begin{longtable}{|
>{\centering\arraybackslash}m{2cm}|
>{\centering\arraybackslash}m{3cm}|
>{\centering\arraybackslash}m{1.5cm}|
}

\caption*{\textbf{Volumi considerati}}
\\

\hline
\textbf{Costrutto} & \textbf{Descrizione} & \textbf{Volume} \\
\hline
\endfirsthead

\hline
\textbf{Costrutto} & \textbf{Descrizione} & \textbf{Volume} \\
\hline
\endhead

Venditore & Numero di venditori & 10 \\
\hline

Acquisto & Numero di vendite al mese & 60 \\
\hline

\end{longtable}

Le operazioni principali individuate sono:

\begin{itemize}
    \item \textbf{Operazione 1 (Scrittura)}: Inserimento di una nuova vendita.
    \item \textbf{Operazione 2 (Lettura)}: Visualizzazione del numero di vendite per ciascun venditore.
\end{itemize}


Di seguito viene effettuata un'analisi dei costi per valutare la convenienza nel mantenere l'attributo \textit{Totale\_Vendite}.

\vspace{0.5cm}
\subsubsection{Scenario con Ridondanza}

\begin{itemize}
    \item{ \textbf{Operazione 1:} Inserimento di una vendita.

       \begin{longtable}{|
            >{\centering\arraybackslash}m{2cm}|
            >{\centering\arraybackslash}m{1.5cm}|
            >{\centering\arraybackslash}m{1cm}|
            }
            
            \hline
            \textbf{Costrutto} & \textbf{Accessi} & \textbf{Tipo} \\
            \hline
            \endfirsthead
            
            \hline
            \textbf{Costrutto} & \textbf{Descrizione} & \textbf{Tipo} \\
            \hline
            \endhead
            
            Venditore & 1 & S \\
            \hline
            
            Acquisto & 1 & S \\
            \hline
            
        \end{longtable}

        Ogni inserimento comporta quindi 2 accessi in scrittura.
        Su base mensile:

        \begin{itemize}
            \renewcommand{\labelitemii}{$\circ$}
            \item 60 vendite → 60 × 2 = 120 accessi in scrittura.
            \item Si considera che una scrittura ha costo doppio rispetto a una lettura. Questo perchè nella maggior parte dei DBMS relazionali la scrittura risulta più "pesante" rispetto alla lettura, dunque si ottiene: → 120 × 2 = 240 unità.
        \end{itemize}
    }
    
    \item {\textbf{Operazione 2}: Lettura mensile.

        \begin{longtable}{|
                >{\centering\arraybackslash}m{2cm}|
                >{\centering\arraybackslash}m{1.5cm}|
                >{\centering\arraybackslash}m{1cm}|
                }
                
                \hline
                \textbf{Costrutto} & \textbf{Accessi} & \textbf{Tipo} \\
                \hline
                \endfirsthead
                
                \hline
                \textbf{Costrutto} & \textbf{Descrizione} & \textbf{Tipo} \\
                \hline
                \endhead
                
                Venditore & 10 & L \\
                \hline
                
        \end{longtable}

        Il per gli accessi in lettura è di 10 unità.
        \newline Si ottiene quindi un \textbf{costo totale mensile} con ridondanza pari a: 240 + 10 = 250 unità.
    }
\end{itemize}

\vspace{0.5cm}
\subsubsection{Scenario senza Ridondanza}

\begin{itemize}
    \item{ \textbf{Operazione 1:} Inserimento di una vendita.

       \begin{longtable}{|
            >{\centering\arraybackslash}m{2cm}|
            >{\centering\arraybackslash}m{1.5cm}|
            >{\centering\arraybackslash}m{1cm}|
            }
            
            \hline
            \textbf{Costrutto} & \textbf{Accessi} & \textbf{Tipo} \\
            \hline
            \endfirsthead
            
            \hline
            \textbf{Costrutto} & \textbf{Descrizione} & \textbf{Tipo} \\
            \hline
            \endhead
               
            Acquisto & 1 & S \\
            \hline
            
        \end{longtable}

        Ogni inserimento comporta 1 solo accesso in scrittura.
        Su base mensile:
        \begin{itemize}
            \renewcommand{\labelitemii}{$\circ$}
            \item 60 vendite → 60 × 1 = 60 accessi in scrittura.
            \item Costo Pesato: → 60 × 2 = 120 unità.
        \end{itemize}
        
    }
    
    \vspace{0.5cm}
    \item {\textbf{Operazione 2}: Lettura mensile.
        \newline Il numero di vendite per venditore può essere ottenuto tramite una query aggregata del tipo:

         \vspace{0.5cm}
        \begin{center}
            \makebox[10cm]{\hspace{2cm}
                SELECT Matricola, COUNT(*) FROM Acquisto GROUP BY Matricola;
            \hspace{2cm}}
        \end{center}
        \vspace{0.5cm}
         
       Anche con un indice su ACQUISTO(Matricola), il sistema deve comunque scorrere tutte le 60 tuple della relazione.
       
        \begin{longtable}{|
                >{\centering\arraybackslash}m{2cm}|
                >{\centering\arraybackslash}m{1.5cm}|
                >{\centering\arraybackslash}m{1cm}|
                }
                
                \hline
                \textbf{Costrutto} & \textbf{Accessi} & \textbf{Tipo} \\
                \hline
                \endfirsthead
                
                \hline
                \textbf{Costrutto} & \textbf{Descrizione} & \textbf{Tipo} \\
                \hline
                \endhead
                
                Acquisto & 60 & L \\
                \hline
                
        \end{longtable}

        Il costo per gli accessi in lettura è dunque di 60 unità.
        \newline Si ottiene quindi un \textbf{costo totale mensile} senza ridondanza pari a: 120 + 60 = 180 unità.
    }
\end{itemize}

\vspace{0.5cm}
\subsubsection{Confronto Finale} 
\vspace{0.5cm}
\begin{longtable}{|
                >{\centering\arraybackslash}m{3.5cm}|
                >{\centering\arraybackslash}m{2.5cm}|
                }
                
                \hline
                \textbf{Scenario} & \textbf{Costo Mensile} \\
                \hline
                \endfirsthead
                
                \hline
                \textbf{Scenario} & \textbf{Costo Mensile} \\
                \hline
                \endhead
                
                Con Ridondanza & 250 \\
                \hline

                Senza Ridondanza & 180 \\                
                \hline
                
        \end{longtable}
        
Dall'analisi emerge che la soluzione senza ridondanza è più efficiente, con un risparmio di circa il 28\%. Questo è dovuto al fatto che l'operazione di lettura aggregata è poco frequente (una sola volta al mese) e il volume di dati è contenuto (60 vendite mensili).

Il costo di aggiornare l'attributo ad ogni inserimento supera quindi il beneficio ottenuto nell'evitare il calcolo durante la fase di lettura.

\subsubsection{Conclusione} 

Si decide pertanto di eliminare l'attributo ridondante Totale\_Vendite dall'entità VENDITORE. Il numero di vendite verrà calcolato direttamente sulla relazione ACQUISTO mediante interrogazioni SQL con COUNT e GROUP BY.

\vspace{0.5cm}

\textbf{Nota:} È importante sottolineare che la scelta dipende dal contesto: qualora la frequenza di lettura fosse significativamente più alta (ad esempio più volte al giorno), il mantenimento della ridondanza potrebbe risultare conveniente. Nel caso in esame, invece, la soluzione ottimale è la sua rimozione.

\vspace{0.5cm}
\subsection{Eliminazione delle Generalizzazioni}

Come evidenziato nello schema E-R in \textbf{Figura 1} le generalizzazioni presenti devono essere eliminate con l'obiettivo di semplificare la successiva implementazione del modello relazionale.
A questo scopo si sviluppa lo schema ristrutturato in \textbf{Figura 2}, dove si aggiunge anche la volontà di mantenere sia la distinzione logica dei dati e sia di evitare valori nulli.
Di seguito vengono quindi elencate le trasformazioni delle due generalizzazioni presenti:

\begin{itemize}
    \item \textbf{Veicolo:} La generalizzazione di \textit{Veicolo} viene eliminata introducendo due nuove relazioni, \textbf{IS\_USATO} e \textbf{IS\_NUOVO}.

    Come già anticipato, questa scelta consente di mantenere una separazione logica tra le due tipologie, evitando la presenza di valori nulli e preservando al contempo lo stesso insieme di attributi già definiti nella fase precedente di ristrutturazione.
    \newline In questo modo, gli identificatori di \textit{Veicolo\_Nuovo} e \textit{Veicolo\_Usato} coincidono rispettivamente con quelli delle relazioni associate.

    \vspace{0.5cm}

    \item \textbf{Dipendente: } In modo analogo, la generalizzazione di \textit{Dipendente} viene eliminata introducendo due nuove relazioni di tipo \textbf{IS\_A}.
    \newline Anche in questo caso, tali relazioni permettono di mantenere la separazione tra \textit{Meccanico} e \textit{Venditore}, evitando al contempo la presenza di valori nulli.
    
\end{itemize}

\vspace{0.5cm}
\subsection{Reificazione della Relazione Ternaria Acquisto}
La relazione ternaria ACQUISTO (tra CLIENTE, VEICOLO, VENDITORE) viene trasformata in un'entità. Dato il vincolo (0,1) dal lato VEICOLO (un veicolo può essere oggetto al più di un acquisto), la chiave primaria dell'entità ACQUISTO è VIN. Le relazioni binarie risultanti sono:

\begin{itemize}
    \item \textbf{EFFETTUA} tra \textit{Cliente} e \textit{Acquisto}: (1,N) e (1,1)
    \item \textbf{OGGETTO\_DI} tra \textit{Veicolo} e \textit{Acquisto}: (0,1) e (1,1)
    \item \textbf{GESTISCE} tra \textit{Venditore} e \textit{Acquisto}: (0,N) e (1,1)
\end{itemize}
L'entità Acquisto possiede gli attributi \texttt{Data\_Acquisto} e \texttt{Prezzo} (entrambi obbligatori).

\vspace{0.5cm}
\subsection{Schema Relazionale}
\begin{figure}[H]
    \begin{adjustwidth}{-2cm}{-2cm}
        \centering
        \label{fig:er-restructured}
        \caption{Diagramma Ristrutturato}
        \includegraphics[width=1.25\textwidth, keepaspectratio]{allegati/Schema Ristrutturato.jpg}
    \end{adjustwidth}
\end{figure}

Il diagramma E-R ristrutturato in \textbf{Figura 2} deriva dallo schema concettuale originario mediante le trasformazioni descritte nelle sezioni precedenti. 
\newline Dopo la ristrutturazione compaiono due nuovi vincoli che non possono essere rappresentati tramite lo schema stesso, ovvero:
\begin{itemize}
    \item Se un \textit{Dipendente} è un \textit{Meccanico}, allora la sua matricola dovrà essere obbligatoriamente presente solo come occorrenza nell'entità \textit{Meccanico}, ma non in quella di \textit{Venditore}, e viceversa.
    \item Se un \textit{Veicolo} è usato, allora il suo VIN dovrà essere obbligatoriamente presente solo come occorrenza nell'entità \textit{Veicolo\_Usato}, ma non in quella di \textit{Veicolo\_Nuovo}, e viceversa.
\end{itemize}
Di seguito si illustra lo schema logico relazionale ottenuto da tale ristrutturazione: 
\begin{itemize}
    \item \textbf{VEICOLO}( \underline{VIN}, Anno\_Immatricolazione, Denominazione\_Commerciale, Codice\_Telaio ) 
    \begin{itemize}
        \item VEICOLO.(Denominazione\_Commerciale, Codice\_Telaio) $\rightarrow$ MODELLO.(Denominazione\_Commerciale, Codice\_Telaio)
    \end{itemize}
    \item \textbf{VEICOLO\_USATO}( \underline{VIN}, Data\_Ultima\_Revisione, Numero\_Proprietari )
    \begin{itemize}
        \item VEICOLO\_USATO.VIN $\rightarrow$ VEICOLO.VIN
    \end{itemize}
    \item \textbf{VEICOLO\_NUOVO}( \underline{VIN}, Garanzia )
    \begin{itemize}
        \item VEICOLO\_NUOVO.VIN $\rightarrow$ VEICOLO.VIN
    \end{itemize}
    \item \textbf{OPTIONAL}( \underline{Codice\_Optional}, Categoria, Nome )
    \item \textbf{MOTORE}(\underline{Codice\_Motore}, Cilindrata, Frazionamento, Alimentazione)
    \item \textbf{MODELLO}( \underline{Denominazione\_Commerciale, Codice\_Telaio}, Anno\_Prima\_Produzione)
    \item \textbf{RICAMBIO}( \underline{OEM}, Nome, Prezzo\_Unitario )
    \item \textbf{INTERVENTO}( \underline{Veicolo, Numero\_Intervento}, Data\_Intervento, Costo, Ore\_Manodopera, Meccanico )
    \begin{itemize}
        \item INTERVENTO.Veicolo $\rightarrow$ VEICOLO.VIN
        \item INTERVENTO.Meccanico $\rightarrow$ MECCANICO.Matricola
    \end{itemize}
    \item \textbf{CLIENTE}( \underline{Codice\_Fiscale}, Nome, Cognome, Data\_Nascita, Telefono )
    \item \textbf{DIPENDENTE}( \underline{Matricola}, Nome, Cognome, Data\_Assunzione, Stipendio )
    \item \textbf{MECCANICO}( \underline{Matricola}, Specializzazione, Livello\_Aziendale )
    \begin{itemize}
        \item MECCANICO.Matricola $\rightarrow$ DIPENDENTE.Matricola
    \end{itemize}
    \item \textbf{VENDITORE}( \underline{Matricola}, Percentuale\_Commissione )
    \begin{itemize}
        \item VENDITORE.Matricola $\rightarrow$ DIPENDENTE.Matricola
    \end{itemize}
    \item \textbf{ACQUISTO}( \underline{Veicolo}, Data\_Acquisto, Cliente, Venditore, Prezzo )
    \begin{itemize}
        \item ACQUISTO.Veicolo $\rightarrow$ VEICOLO.VIN
        \item ACQUISTO.Cliente $\rightarrow$ CLIENTE.Codice\_Fiscale
        \item ACQUISTO.Venditore $\rightarrow$ VENDITORE.Matricola
    \end{itemize}
\end{itemize}

\vspace{0.5cm}
Di seguito si elencano le tabelle ottenute dalle relazioni molti a molti: 
\begin{itemize}
    \item \textbf{EQUIPAGGIA}( \underline{Optional, Veicolo} )
    \begin{itemize}
        \item EQUIPAGGIA.Optional $\rightarrow$ OPTIONAL.Codice\_Optional
        \item EQUIPAGGIA.Veicolo $\rightarrow$ VEICOLO\_NUOVO.VIN
    \end{itemize}
    \item \textbf{UTILIZZA}( \underline{Ricambio, Numero\_Intervento, Veicolo}, Quantità )
    \begin{itemize}
        \item UTILIZZA.Ricambio $\rightarrow$ RICAMBIO.OEM
        \item UTILIZZA.Numero\_Intervento $\rightarrow$ INTERVENTO.Numero\_Intervento
        \item UTILIZZA.Veicolo $\rightarrow$ INTERVENTO.Veicolo
    \end{itemize}
    \item \textbf{MOTORIZZATO\_DA} ( \underline{Codice\_Telaio, Denominazione\_Commerciale, Motore}, Potenza )
    \begin{itemize}
        \item MOTORIZZATO\_DA.(Codice\_Telaio, Denominazione\_Commerciale) $\rightarrow$ MODELLO.(Codice\_Telaio, Denominazione\_Commerciale)
        \item MOTORIZZATO\_DA.Motore $\rightarrow$ MOTORE.Codice\_Motore
    \end{itemize}
\end{itemize}

